\apartado{Identidad, IAM y almacenamiento}
\subsection{Dos niveles de identidad}
Firebase Authentication identifica al estudiante; el backend debe validar el ID token antes de confiar en su usuario \cite{firebase-token}. La autorización de producto añade la lista de cuentas habilitadas y comprueba propietario o concesión en cada operación. Que una persona haya iniciado sesión no le permite leer una materia ajena.

IAM controla identidades de infraestructura. Se proponen cuentas separadas para API, worker, despacho de tareas y despliegue, con permisos mínimos sobre los recursos que necesitan. No se incluyen llaves de servicio en el repositorio. Los secretos del proveedor de IA se reservarían al servidor mediante Secret Manager.

Las bibliotecas servidor de Firestore se autorizan mediante IAM; no sustituyen la autorización del dominio \cite{firestore-iam}. La propuesta deniega accesos directos del cliente a colecciones y objetos. Por ello, una regla cliente restrictiva no basta para proteger una ruta de servidor que omita comprobar el propietario.

\subsection{Base documental frente a almacenamiento de objetos}
\begin{tabularx}{\linewidth}{P{2.4cm}Y Y}
\toprule
\textbf{Dimensión} & \textbf{Firestore: documentos} & \textbf{Storage: objetos}\\
\midrule
Contenido & Perfil, materia, tema, agenda, revisión y estado de trabajo. & PDF, PPTX, imágenes y paquetes derivados.\\
Uso esperado & Consultas acotadas e índices; cambios pequeños por entidad. & Guardar y recuperar bytes por clave y versión.\\
Organización & Alcance por usuario y relaciones verificadas por servidor. & Claves privadas calculadas por servicio; no por el navegador.\\
Consumo a medir & Lecturas, escrituras, almacenamiento e índices. & Bytes, operaciones, transferencias y retención.\\
Consistencia & \codigo{revision} y \codigo{expectedRevision} para detectar conflictos. & Generación de objeto y versión de extracción para evitar confusión.\\
\bottomrule
\end{tabularx}

El término ``documento'' en Firestore se refiere a un registro estructurado; no significa almacenar allí el PDF completo. Separar metadatos y archivos permite consultar la organización sin descargar material pesado, y aplicar políticas específicas de tamaño, versión y eliminación.

\subsection{Carga privada prevista para etapas posteriores}
La API resolvería una clave de objeto autorizada y emitiría un ticket de corta duración. Una URL firmada concede acceso temporal a quien la posea \cite{signed-url}; no se registraría en logs ni se trataría como enlace público de biblioteca. Completar la carga exigiría verificar objeto, tamaño y tipo real. La disponibilidad de este diseño no significa que el flujo de upload ya esté integrado en H1.

Los originales, extracciones y recursos derivados tienen ciclos de vida diferentes. El borrado debe considerar todos sus descendientes y la retención física del proveedor; archivar una materia conserva historial y no equivale a eliminarlo.
